\documentclass[11pt]{article}
\usepackage{acl}
\usepackage{times}
\usepackage{latexsym}
\usepackage{amsmath}
\usepackage{booktabs}
\usepackage{graphicx}
\usepackage{microtype}

\title{Cross-Lingual Hallucination Drift in LLMs:\\Does It Depend on Task Type?\\
\normalsize Midway Report}

\author{
  Bhoomika Monthy Rajashekar \quad Devinn Chi \quad Chun Hsu \quad Anagha P Krishna \\
  Boston University \\
  \texttt{\{bhoomi02, dhchi, rexhsu, anaghapk\}@bu.edu} \\
  Advisor: Aaron Mueller
}

\begin{document}
\maketitle

%% ABSTRACT
\begin{abstract}
Large Language Models (LLMs) are known to hallucinate in non-English languages,
but whether this cross-lingual drift is task-dependent remains an open question.
We test this by measuring hallucination rates of Aya Expanse 8B across English,
Spanish, and Swahili on two structurally different tasks: factual QA (TruthfulQA)
and commonsense reasoning (XCOPA), using GPT-4o-mini as an LLM-as-a-Judge evaluator.
Our completed experiments yield a striking result: hallucination drift is strongly
task-dependent. On TruthfulQA, Spanish drift is negligible ($\Delta$HR $= {-2.66}$~pp,
$p = 0.60$). On XCOPA, Swahili drift is catastrophic ($\Delta$HR $= {+90.67}$~pp,
$p < 0.001$). The cross-task drift interaction score $\Phi = {-93.33}$~pp is
statistically significant ($\chi^2 = 170.62$, $p < 0.001$). A secondary finding
is that when hallucinating on XCOPA in Swahili, the model generates nearly twice
as many tokens as on faithful responses (186 vs.\ 100), suggesting incoherent
rambling rather than confident fabrication. These results extend and refine prior
work by showing that task type is a major confound in cross-lingual hallucination
estimation.
\end{abstract}

%% 1. INTRODUCTION
\section{Introduction}

Hallucination in LLMs---generating content that is factually incorrect or
unfaithful---is a well-documented failure mode \cite{alansari2025survey}.
A natural concern as LLMs are deployed globally is whether hallucination
rates are higher in non-English languages, particularly low-resource ones
where training data is sparse \cite{blasi-etal-2022-systematic}.
This phenomenon is termed \textit{Cross-Lingual Hallucination Drift}: the
change in hallucination rate when a model operates in a non-English language.

A pivotal study by \citet{ulislam-etal-2025-hallucinate} measured token-level
hallucination rates across 30 languages and 6 open-source LLM families, finding
that \textit{length-normalized hallucination rates are uncorrelated with language
resource level}---a surprising result that overturns the common assumption that
low-resource languages hallucinate more. However, that study evaluated a single
task type: open-domain, knowledge-intensive long-form QA. Whether this finding
holds across structurally different task types remains unexplored.

We investigate this gap. Factual recall (TruthfulQA) requires the model to retrieve
internally stored parametric knowledge, while commonsense reasoning (XCOPA) requires
applying world knowledge to select a plausible cause or effect \cite{qi-etal-2023-cross}.
These two task types may expose fundamentally different failure modes across languages.
We formally define the \textbf{Drift Interaction Score}:
\[
\Phi_l = \Delta\mathrm{HR}_{l,\,\text{TruthfulQA}} - \Delta\mathrm{HR}_{l,\,\text{XCOPA}}
\]
where $\Delta\mathrm{HR}_{l,t} = \mathrm{HR}_{l,t} - \mathrm{HR}_{\text{en},t}$.
A large $|\Phi_l|$ indicates drift is task-dependent for language $l$. Our
completed experiments confirm this hypothesis decisively.

%% 2. METHODS
\section{Methods}

\subsection{Datasets}

We use two benchmarks representing structurally different task types.

\paragraph{TruthfulQA} \cite{lin-etal-2022-truthfulqa} consists of 817 questions
across 38 domains designed to elicit confident false beliefs. We use the English
original and the Spanish translation from the \texttt{alexandrainst/m\_truthfulqa}
HuggingFace dataset. Spanish is a mid-resource language with substantial LLM
training data coverage.

\paragraph{XCOPA} is a cross-lingual commonsense causal reasoning benchmark
spanning 11 typologically diverse languages. Each example presents a premise and
two candidate causes or effects; the model must select the more plausible one. We
use the English split from SuperGLUE COPA and the Swahili (\texttt{sw}) split.
Swahili is a low-resource language. Spanish is not available in XCOPA, so our
cross-task comparison uses Spanish for TruthfulQA and Swahili for XCOPA.

We sample 150 examples per language per task (random seed 42), yielding 600 total
evaluation instances across four language--task cells.

\subsection{Models}

\paragraph{Target model.} We evaluate \textbf{Aya Expanse 8B}, a multilingual
open-source instruction-tuned model by Cohere Labs, queried with zero-shot prompts
written in the target language for each cell.
Prompts are templated separately for TruthfulQA (free-form answer) and XCOPA
(choose A or B).

\paragraph{Judge model.} We use \textbf{GPT-4o-mini} as our LLM-as-a-Judge
evaluator, following the methodology of \citet{kang-etal-2024-comparing}. The judge
receives each (question, model response) pair and labels it \textit{Hallucinated} or
\textit{Faithful} with a brief justification. We chose GPT-4o-mini for its strong
multilingual comprehension across all three languages.

\subsection{Evaluation Pipeline}

The pipeline proceeds in five stages: (1) dataset download and sampling; (2) zero-shot
inference with Aya Expanse 8B on all four language--task cells; (3) GPT-4o-mini judge
labeling, with automatic retry for any \texttt{ERROR} labels; (4) hallucination rate
computation; and (5) statistical significance testing.

We compute the hallucination rate as:
\[
\mathrm{HR}_{l,t} = \frac{\#\,\text{Hallucinated}}{N_{l,t}} \times 100\,(\%)
\]
where $N_{l,t} = 150$ for all cells. All code is publicly available in our
repository.

%% 3. EXPERIMENTS
\section{Experiments}

\subsection{Setup}

All experiments run on Python 3.10 with \texttt{pandas}, \texttt{scipy},
\texttt{matplotlib}, and the HuggingFace \texttt{datasets} library. Aya Expanse 8B
inference requires a GPU; judge calls use the OpenAI API. Total labeled examples: 600
(150 per cell). Judge labels were inspected and retried for any malformed outputs.

\subsection{Statistical Tests}

We apply three tests to assess significance of hallucination differences:

\begin{itemize}
  \item \textbf{Mann-Whitney U} (per task): tests whether hallucination rates differ
        between English and the non-English language within a task.
  \item \textbf{Chi-squared} ($2{\times}2$ contingency): tests whether hallucination
        outcome is independent of task type for non-English responses---directly testing
        the $\Phi$ hypothesis.
  \item \textbf{Fisher's exact}: robustness check on the same $2{\times}2$ table,
        preferred for smaller cell counts.
\end{itemize}

%% 4. RESULTS
\section{Results}

\subsection{Hallucination Rates and Drift}

Table~\ref{tab:hr} shows hallucination rates across all four language--task cells.
TruthfulQA rates are comparable between English (27.33\%) and Spanish (24.67\%),
yielding a negligible and non-significant drift of $-2.66$~pp. In contrast, XCOPA
Swahili hallucination is near-total (98.67\%), representing a drift of $+90.67$~pp
above the English XCOPA baseline (8.00\%).

\begin{table}[t]
\centering
\small
\begin{tabular}{llrr}
\toprule
\textbf{Task} & \textbf{Language} & \textbf{HR (\%)} & \textbf{$\Delta$HR (pp)} \\
\midrule
TruthfulQA & English  & 27.33 & ---      \\
TruthfulQA & Spanish  & 24.67 & $-2.66$  \\
XCOPA      & English  &  8.00 & ---      \\
XCOPA      & Swahili  & 98.67 & $+90.67$ \\
\bottomrule
\end{tabular}
\caption{Hallucination rates (HR) and drift ($\Delta$HR vs.\ English baseline)
per language--task cell. $N=150$ per cell.}
\label{tab:hr}
\end{table}

\subsection{Drift Interaction and Statistical Significance}

Table~\ref{tab:stats} presents all statistical test results. The XCOPA en$\to$sw
drift is highly significant ($p < 0.001$), while TruthfulQA en$\to$es drift is not
($p = 0.60$). The cross-task $\chi^2$ test confirms that hallucination outcome is not
independent of task type ($\chi^2 = 170.62$, $p < 0.001$), validated by Fisher's
exact test (OR $= 226.0$, $p < 0.001$). The aggregate drift interaction score is
$\Phi = {-93.33}$~pp, far exceeding our pre-specified minimum effect threshold of
$|\Phi| \geq 2$~pp.

\begin{table*}[t]
\centering
\small
\begin{tabular}{llrrl}
\toprule
\textbf{Test} & \textbf{Comparison} & \textbf{Stat.} & \textbf{$p$} & \textbf{Sig.} \\
\midrule
Mann-Whitney U & TruthfulQA: en vs es & 11550.0  & 0.600    & No  \\
Mann-Whitney U & XCOPA: en vs sw      & 1050.0   & $<$0.001 & Yes \\
Chi-squared    & es/TQA vs sw/XCOPA   & 170.62   & $<$0.001 & Yes \\
Fisher's exact & es/TQA vs sw/XCOPA   & OR=226.0 & $<$0.001 & Yes \\
\bottomrule
\end{tabular}
\caption{Statistical test results. TQA = TruthfulQA. Chi-squared and Fisher's
exact test the $\Phi$ hypothesis directly.}
\label{tab:stats}
\end{table*}

\subsection{Token Verbosity}

A secondary finding: when hallucinating on XCOPA in Swahili, Aya generates an average
of 186 tokens per response, compared to 100 tokens for faithful responses in the same
cell---an 86\% increase. This pattern does not hold for TruthfulQA in either language,
suggesting that XCOPA Swahili hallucinations reflect \textit{incoherent rambling}
rather than the confident but wrong short answers typical of factual QA hallucination.

\begin{figure}[h]
  \centering
  \includegraphics[width=\columnwidth]{hr_by_language_task.png}
  \caption{Hallucination rates by language and task. XCOPA Swahili is near-total
  (98.67\%), while TruthfulQA drift between English and Spanish is negligible.}
  \label{fig:hr}
\end{figure}

\begin{figure}[h]
  \centering
  \includegraphics[width=\columnwidth]{drift_by_task.png}
  \caption{Cross-lingual drift ($\Delta$HR vs.\ English) by language--task cell.
  Spanish TruthfulQA shows near-zero drift ($-2.7$~pp), while Swahili XCOPA
  shows catastrophic drift ($+90.7$~pp).}
  \label{fig:drift}
\end{figure}

\subsection{Discussion}

Our results provide strong evidence that cross-lingual hallucination drift is
task-dependent. The contrast between tasks is stark: a multilingual model that
handles Spanish factual QA nearly as well as English collapses almost entirely on
Swahili commonsense reasoning. This suggests that commonsense reasoning in XCOPA
may be more sensitive to linguistic and cultural grounding in the target language---the
model cannot fall back on parametric factual recall and instead produces incoherent
output. This finding directly extends \citet{ulislam-etal-2025-hallucinate}: their
null correlation between HR and language resource level may hold for knowledge-intensive
QA but not for commonsense reasoning.

One limitation is that Spanish and Swahili appear in different tasks (TruthfulQA
and XCOPA, respectively), since Spanish is not available in XCOPA. Our $\Phi$ score
therefore conflates task-type and language-identity differences. We address this by
framing $\Phi$ as a cross-task aggregate comparison rather than a per-language score,
and note it as a priority for the final report.

%% 5. PROJECT PLAN
\section{Project Plan}

All core experiments are complete. The remaining work for the final report is:

\begin{enumerate}
  \item \textbf{Qualitative error analysis.} Categorize GPT-4o-mini judge reasons
        (e.g., incoherent, wrong answer, fabricated, code-switching) and report
        hallucination type distributions per cell.
  \item \textbf{XCOPA question-type breakdown.} Separate XCOPA hallucination rates
        by question type (\textit{cause} vs.\ \textit{effect}) to assess whether
        causal direction affects the failure mode.
  \item \textbf{Dual-judge agreement.} Run a second judge on a 50-sample subset
        per cell to compute inter-annotator agreement (Cohen's $\kappa$) and validate
        GPT-4o-mini reliability on Swahili.
  \item \textbf{Final report write-up.} Expand the methods and discussion sections,
        incorporate error analysis findings, and polish figures for submission.
\end{enumerate}

%% TEAM CONTRIBUTIONS
\clearpage
\section*{Team Contributions}

\begin{table}[h]
\centering
\small
\begin{tabular}{p{2.2cm} p{1.8cm} p{4.2cm}}
\toprule
\textbf{Member} & \textbf{Role} & \textbf{Contributions} \\
\midrule
Chun Hsu
  & Data \& Infra
  & Dataset download and preprocessing;
    HuggingFace pipeline setup;
    prompt formatting and JSON storage \\
\addlinespace
Devinn Chi
  & Experiments
  & Zero-shot inference with Aya Expanse 8B
    across all language--task cells;
    GPT-4o-mini judge runs;
    error label retry pipeline \\
\addlinespace
Bhoomika Monthy Rajashekar
  & Evaluation \& Metrics
  & LLM-as-a-Judge prompt design;
    hallucination rate computation;
    $\Delta$HR and $\Phi$ metric implementation;
    Streamlit dashboard \\
\addlinespace
Anagha P Krishna
  & Analysis \& Writing
  & Statistical tests (Mann-Whitney,
    chi-squared, Fisher's exact);
    visualizations; report write-up \\
\bottomrule
\end{tabular}
\caption{Team task distribution for the midway report.}
\label{tab:team}
\end{table}

%% REFERENCES
\bibliographystyle{acl_natbib}
\bibliography{references}

\end{document}
